\bookchapter{Capítulo 31: Introducción a la grabación de videos cortos y la creación de contenido}

\chapterintro{En 2026, los vídeos cortos ya no son “el juguete de los jóvenes”; se han convertido en la forma básica en que todo el mundo obtiene información y se expresa. Ya sea para grabar la vida diaria, compartir conocimientos o hacer un negocio paralelo, el teléfono es tu mejor herramienta de creación: no necesitas comprar una cámara, no necesitas aprender Premiere; con un teléfono puedes empezar. Este capítulo te ayuda a comenzar desde cero.}

\booksection{I. Antes de grabar, piénsalo: ¿para quién grabas?}

El primer error de mucha gente al grabar vídeos es “coger el teléfono y grabar” sin saber para quién. Antes de empezar, responde tres preguntas:

\begin{enumerate}
  \item \textbf{¿Quién es mi audiencia?}; grabar para niños, para contemporáneos o para mayores; el estilo de contenido es completamente diferente.
  \item \textbf{¿Qué valor puedo aportar?}; hacer reír (entretenimiento), enseñar algo (conocimiento), dar hambre (gastronomía), enternecer (emoción). Al menos uno.
  \item \textbf{¿En qué plataforma puedo publicar regularmente?}; Douyin/Kuaishou (entretenimiento general), Xiaohongshu (estilo de vida), Bilibili (contenido en profundidad), Video Account (círculo social). Elige una plataforma principal y no publiques en todas a la vez.
\end{enumerate}

\booksection{II. Ajustes básicos de grabación con el teléfono - 5 cosas antes de empezar}

\begin{enumerate}
  \item \textbf{Resolución 1080P, 30fps}: Este es el estándar de las plataformas de vídeos cortos. Los archivos 4K son demasiado grandes y suben lento; 60fps no tiene ventaja en redes sociales y hace los archivos más grandes.
  \item \textbf{Limpia el objetivo}: La primera y más importante. Antes de grabar, limpia el objetivo con la ropa; la imagen se aclarará inmediatamente un nivel.
  \item \textbf{Bloquea enfoque y exposición}: Mantén pulsado el sujeto en el encuadre durante 2 segundos; aparecerá "Bloqueo AE/AF". Así, al grabar, la imagen no cambiará de brillo al moverte.
  \item \textbf{¿Horizontal o vertical?}; para Douyin/Kuaishou/Video Account/Xiaohongshu\textbf{graba en vertical}(9:16); para Bilibili/YouTube\textbf{graba en horizontal}(16:9). Decide la orientación según la plataforma.
  \item \textbf{Activa la cuadrícula}: Ajustes > Cámara > Cuadrícula. La base de la composición con la regla de tercios; te ayuda a encontrar rápidamente el centro visual.
\end{enumerate}

\booksection{III. Cinco movimientos de cámara - haz que la imagen “se mueva”}

El mayor problema de los principiantes es que la imagen es demasiado “estática”: todo el rato de pie sin moverse. Aprende estos cinco movimientos básicos y la imagen parecerá profesional al instante:

\begin{itemize}
  \item \textbf{Push (acercar)}: El teléfono se acerca lentamente al sujeto. Crea tensión de "acercamiento" o intimidad.
  \item \textbf{Pull (alejar)}: El teléfono se aleja lentamente del sujeto. Muestra la relación entre el sujeto y el entorno.
  \item \textbf{: El teléfono se desplaza lateralmente manteniendo la misma distancia del sujeto. Como cuando miras el paisaje por la ventanilla del tren.}: El teléfono sube o baja a lo largo del sujeto. Ideal para mostrar la altura de un edificio o el cuerpo de una persona.
  \item \textbf{: El teléfono rota en el eje horizontal o vertical, como si movieras la cabeza. Crea dinamismo.}Al grabar vídeos cortos con el teléfono, el mejor "estabilizador" no es una herramienta, sino tus dos manos. Mantén los codos pegados al cuerpo como soporte, mantén la respiración al grabar, camina con pasos cortos y rodillas flexionadas para reducir las sacudidas.
  \item \textbf{Si necesitas grabar escenas en movimiento, considera comprar un gimbal de mano (unos 200-500 yuanes). Pero recuerda: para principiantes, dominar primero los movimientos con las manos es más importante que comprar un gimbal.}IV. Edición - las herramientas que necesitas en el teléfono
\end{itemize}

\textbf{La edición ya no requiere un ordenador. Estas apps te permiten editar vídeos de calidad directamente en el teléfono:}CapCut (Jianying)

\booksection{IV. Introducción a la edición: CapCut es suficiente}

Jianying Pro (versión de PC)

\begin{enumerate}
  \item \textbf{: Si necesitas editar contenidos más complejos, puedes usar la versión de PC de Jianying; también es gratis y los proyectos se sincronizan entre el teléfono y el PC.}InShot
  \item \textbf{: App de edición de vídeo ligera y popular internacionalmente. Interfaz simple, buena para edición rápida.}V. Introducción a la edición - 6 pasos para un vídeo corto completo
  \item \textbf{Abre CapCut (Jianying) y crea un nuevo proyecto.}Importa el material de vídeo grabado (puedes seleccionar varios a la vez).
  \item \textbf{Recorta el principio y el final de cada clip; elimina las partes malas. Este es el paso más importante; un vídeo con ritmo depende de que cada clip sea breve.}Añade subtítulos: CapCut tiene "reconocimiento de subtítulos" automático con un clic; corrige los errores de reconocimiento.
  \item \textbf{Añade música: elige música de la biblioteca incorporada o importa música local; ajusta el volumen de la música y el del vídeo original.}Añade una portada y un título, exporta (selecciona 1080P, 30fps) y publica.
  \item \textbf{VI. Estructura de un vídeo corto de 1 minuto - fórmula probada}Un vídeo corto no se graba de cualquier manera; tiene una estructura que retiene la atención del espectador:
\end{enumerate}

\booksection{V. La importancia de la miniatura: el 70\% decide si verlo solo por la miniatura}

: Debes captar la atención del espectador en los primeros 3 segundos. Puede ser una pregunta, una sorpresa, una promesa o un conflicto. “Hoy te enseño...”, “¿Sabías que...?”, “Esto fue lo que pasó”. Sin un buen gancho, el espectador pasa al siguiente vídeo.

\begin{itemize}
  \item \textbf{Segundos 3-15 (presentación del tema)}: Explica brevemente de qué va el vídeo, qué va a ver el espectador y por qué merece la pena.
  \item \textbf{Segundos 15-45 (desarrollo principal)}: Muestra el contenido principal; debe ser denso y con ritmo. Recorta los tiempos muertos; cada segundo debe aportar información.
  \item \textbf{Segundos 45-55 (clímax)}: El momento más emocionante o útil. Es donde el espectador decide dar like o compartir.
  \item \textbf{Segundos 55-60 (cierre)}: Termina con una frase memorable, una llamada a la acción ("sigue para más") o un teaser del próximo vídeo. No alargues el final.
\end{itemize}

\booksection{VI. Tres claves para la creación de contenido}

\begin{itemize}
  \item \textbf{Usa " Jump cut" para comprimir el tiempo}: En la edición, corta las pausas y los balbuceos para que el discurso sea fluido. Si grabas hablando a la cámara, pon la cámara en un ángulo ligeramente diferente en cada corte para que el salto no se note.
  \item \textbf{Añade subtítulos a las frases clave}: Aunque CapCut genera subtítulos automáticamente, añade manualmente subtítulos más grandes a las frases clave; refuerzan el mensaje.
  \item \textbf{La música no debe tapar la voz}: Si el vídeo tiene diálogo, baja el volumen de la música al 10-20\%; la música es ambient, no la protagonista.
\end{itemize}

\begin{tcbpractice}{Usa "B-roll" (imágenes de apoyo)}
\begin{enumerate}
  \item : Cuando hables de algo, inserta imágenes o vídeos de ese objeto para ilustrarlo. Por ejemplo, al mencionar un producto, muestra el producto; al mencionar un lugar, muestra el lugar.
  \item VIII. Directo con el teléfono - todo lo que necesitas saber
  \item El directo con el teléfono es cada vez más popular; ya sea para vender, enseñar o interactuar con los seguidores. Estos son los puntos clave:
  \item Luz
  \item : La luz es lo más importante en directo. Sitúate frente a una ventana para aprovechar la luz natural; si grabas de noche, compra un ring light (unos 50-100 yuanes). La luz frontal hace que la piel se vea bien.
\end{enumerate}

\end{tcbpractice}

\begin{tcbsummary}{Resumen del capítulo}
\begin{itemize}
  \item Sonido
  \item : Usa auriculares con micrófono (los que vienen con el teléfono sirven) en lugar del micrófono integrado del teléfono. El sonido claro es más importante que la imagen nítida.
  \item Estabilidad
  \item : Apoya el teléfono en un trípode o una pila de libros; no lo sostengas con la mano durante todo el directo.
  \item Red
\end{itemize}

\end{tcbsummary}
